\chapter{Конструкторский раздел}

В данном разделе представлена архитектура программно-алгоритмического комплекса,
описаны его компоненты, их взаимодействие и ключевые структуры данных, приведены 
реализуемые схемы алгоритмов, описаны сущности и ограничения целостности проектируемой базы данных.

\section{Архитектура разрабатываемого программно-алгоритмического комплекса}

Архитектура комплекса будет представлена моделью клиент-сервер, где клиент ---
приложение, позволяющее пользователю управлять списком доверенных устройств, а сервер
--- приложение, отвечающее непосредственно за монтирование устройств, хранение 
записей о монтировании, работе с базой данных. Схематично данная модель представлена
на рисунке~\ref{img:client_server.drawio}.

\includeimage
    {client_server.drawio}
    {f}
    {H}
    {0.80\textwidth}
    {Модель клиент-сервер}

\section{Взаимодействие компонентов программно-алгоритмического комплекса}

Взаимодействие клиента и сервера может быть описано с помощью спецификации OpenAPI~\cite{OpenAPI}. 
Использование OpenAPI позволяет формализовать структуру запросов и ответов, определить 
параметры запросов и форматы передаваемых данных. 

На рисунке~\ref{img:api} представлено описание программного интерфейса приложения, 
включающее маршруты взаимодействия клиента и сервера.

\includeimage
    {api}
    {f}
    {H}
    {1\textwidth}
    {Описание программного интерфейса приложения}

\section{Ключевые структуры данных}

Серверная часть использует несколько структур данных для отправки ответов клиенту и 
использования во внутренних методах.

\begin{itemize}[label=---]
    \item DeviceInfo --- структура формального описания USB-устройства:
        \begin{enumerate}
        \item vendorId --- идентификатор производителя;
        \item productId --- идентификатор устройства;
        \item serial --- серийный номер;
        \item vendorName --- название производителя;
        \item productName --- название продукта.
        \end{enumerate}
    \item Device --- структура описания доверенного устройства, используется для передча на клиент при запросе списка
    доверенных устройств:
        \begin{enumerate}
        \item id --- первичный ключ из базы данных;
        \item info --- формальное описание DeviceInfo;
        \item validTo --- дата истечения доверия.
        \end{enumerate}
    \item EventType --- перечисляемый тип, используемый для обозначения физического манипулирования
    с накопителем, где INSERT(0) --- подключение устройства, REMOVE(1) --- извлечение устройства;
    \item DeviceEvent --- структура, описывающая физическое событие с накопителем. Имеет поля type 
    типа EventType, и devNode --- узел устройства;
    \item MODE --- перечисляемый тип, используемый для обозначения режима монтирования, где 
    RO(0) --- режим только для чтения, а RW(1) --- режим полного доступа;
    \item MountRecord --- структура описания события монтирования, используется для передачи на клиент при запросе подключённых
    в данный момент устройств:
        \begin{enumerate}
        \item id --- первичный ключ устройства из базы данных(при наличии);
        \item devNode --- узел устройства;
        \item mountPoint --- путь монтирования;
        \item info --- формальное описание DeviceInfo;
        \item mode --- режим монтирования.
        \end{enumerate}
\end{itemize}

\section{Схема алгоритма монтирования устройства}

На рисунке~\ref{img:mount_algo.drawio} представлена схема алгоритма монтирования устройства
при его подключении к компьютеру.

\includeimage
    {mount_algo.drawio}
    {f}
    {H}
    {0.50\textwidth}
    {Схема алгоритма монтирования}

\section{Схема алгоритма удаления устройства из списка доверенных устройств}

На рисунке~\ref{img:delete_from_white_list_algo.drawio} представлена схема алгоритма 
удаления устройства из списка доверенных устройств.

\includeimage
    {delete_from_white_list_algo.drawio}
    {f}
    {H}
    {0.45\textwidth}
    {Схема алгоритма удаления устройства}

\section{Схема алгоритма обработки узлов устройств и записей о точках монтирования при инициализации приложения}

На рисунке~\ref{img:bypass_devNodes.drawio} представлена схема алгоритма 
обработки узлов устройств и записей о точках монтирования при инициализации приложения.

\includeimage
    {bypass_devNodes.drawio}
    {f}
    {H}
    {0.50\textwidth}
    {Схема алгоритма обработки узлов устройств и записей о точках монтирования при инициализации приложения}

\section{Схема алгоритма проверки записи о монтировании на предмет актуальности}

На рисунке~\ref{img:refresh_mount_after_reboot_daemon_algo.drawio} представлена схема алгоритма 
проверки записи о монтировании на предмет актуальности.

\includeimage
    {refresh_mount_after_reboot_daemon_algo.drawio}
    {f}
    {H}
    {0.70\textwidth}
    {Схема алгоритма проверки записи о монтировании на предмет актуальности}

\section{Описание сущностей проектируемой базы данных}

База данных содержит следующие таблицы:

\begin{itemize}[label=---]
    \item DeviceInfo --- описание устройства:
    \begin{enumerate}
        \item id --- первичный ключ;
        \item vendorId --- идентификатор производителя;
        \item productId --- идентификатор продукта;
        \item serial --- серийный номер;
        \item productName --- название производителя;
        \item vendorName --- название устройства.
    \end{enumerate}
    \item Device --- описание записи в списке доверенных устройств:
    \begin{enumerate}
        \item id --- первичный ключ;
        \item deviceInfoId --- идентификатор описания устройства;
        \item validTo --- дата окончания периода доверия.
    \end{enumerate}
    \item MountRecord --- описание события монтирования:
    \begin{enumerate}
        \item id --- первичный ключ;
        \item deviceInfoId --- идентификатор описания устройства;
        \item devNode --- узел устройства;
        \item mountPoint --- точка монтирования;
        \item mode --- режим монтирования.
    \end{enumerate}
\end{itemize}

\section{Диаграмма проектируемой базы данных}

На рисунке~\ref{img:database} представлена диаграмма проектируемой базы данных.

\includeimage
    {database}
    {f}
    {H}
    {0.60\textwidth}
    {Диаграмма базы данных}

Между таблицами базы данных образуется связь один к одному.

\section{Описание проектируемых ограничений целостности базы данных}

На значения полей базы данных накладываются следующие ограничения целостности:

\begin{itemize}[label=---]
    \item DeviceInfo:
    \begin{enumerate}
        \item поля vendorId и productId должны содержать только цифры или строчные или 
        заглавные латинские буквы, их длина должна равняться 4-ём символам, они не могут быть пустыми;
        \item поле serial не может быть пустым;
        \item комбинация полей vendorId, productId, serial должна быть уникальной.
    \end{enumerate}
    \item MountRecord:
    \begin{enumerate}
        \item поле deviceInfoId является внешним ключом, ссылающимся на таблицу DeviceInfo, должно быть непустым и уникальным;
        \item devNode должно быть непустым и уникальным;
        \item mountPoint должно быть непустым и уникальным;
        \item mode должно быть непустым и принимать только значения <<RO>> и <<RW>>.
    \end{enumerate}
    \item Device: поле deviceInfoId является внешним ключом, ссылающимся на таблицу DeviceInfo, должно быть непустым и 
    уникальным, а validTo должно содержать дату в формате <<ГГГГ-ММ-ДД>>.
\end{itemize}